\documentclass[10pt,a4paper]{amsart}
\usepackage[margin=19mm]{geometry}
\usepackage{amsmath,amssymb,mathtools}
\usepackage[scr=rsfs,cal=euler]{mathalfa}
\usepackage{xfrac}
\usepackage[unicode,hidelinks,bookmarks=false]{hyperref}
\newcommand{\ie}{i.e.\ }
\newcommand{\cf}{cf.\ }
\newcommand{\resp}{respectively\ }
\newcommand{\Subsection}{\subsection}
\providecommand{\nobreakdash}{\nobreak\hbox{-}}
\setlength{\emergencystretch}{2em}
\multlinegap0pt
\usepackage{xfrac}
\usepackage{mathtools}
\usepackage{amssymb}
\usepackage[shortlabels]{enumitem}
\usepackage[normalem]{ulem}
\newtheorem{lemma}{Lemma}
\newtheorem{theorem}{Theorem}
\newcommand{\Cl}{\textup{Cl}}
\title{(Locally) Associated Subrings in Polynomial and Power Series Extensions}
\author{Grant Moles, Joseph Swanson}
\date{Source: arXiv:2608.07741v1 (CC BY 4.0)}
\begin{document}
\maketitle

\noindent\emph{Source and licence.} (Locally) Associated Subrings in Polynomial and Power Series Extensions. Version arXiv:2608.07741v1 (CC BY 4.0), distributed by arXiv under the Creative Commons Attribution 4.0 International licence, http://creativecommons.org/licenses/by/4.0/, as recorded in the arXiv licence metadata for this exact version and shown on the abstract page https://arxiv.org/abs/2608.07741v1. Full licence notice: \texttt{licenses/arxiv-2608.07741-CC-BY-4.0.txt}.\par\smallskip

\noindent\textit{Editorial scope.} Counted unit: the article's theorem sequence (source0.tex lines 239--251). The article's complete original proof of it is delivered with it (source0.tex lines 253--292, one \texttt{proof} environment of 40 lines). No mathematics is written, repaired or re-derived: the statement and the proof are the article's own lines, in the article's own notation, and no retained line is substituted or re-flowed.  Retained uncounted as the source-local context the proof needs: lines 213--216; lines 226--229.  Nothing was omitted from any retained span, so the package carries no omission record.  No retained line contains a citation, so the wrapper carries no bibliography and no bibliography entry.  No retained line contains a figure environment, an image-inclusion command or drawing code, so no image is delivered and the package declares no asset dependency.  Editorial formatting only, in addition: an AMS \texttt{amsart} preamble that restates the article's own declarations and macros used by the retained lines, namely the environments \texttt{lemma}, \texttt{theorem} on separate counters, together with the standard packages those lines need.  The retained environments print in this document as Lemma 1, Theorem 1 in order of appearance, whereas the article prints the same items with the numbering of its own full document.  The complete native source of the article as distributed in the arXiv e-print for this exact version is bundled unchanged as \texttt{sources/207105/source0.tex}.
\begin{lemma}
    \label{going up or going down}
    Let $T$ be an integral domain, $R$ a subring with identity, and $I=(R:T)$. If $J$ is an ideal of $R$ which is comaximal to $I$, then $JT\cap R=J$. If $A$ is an ideal of $T$ that is comaximal to $I$, then $(R\cap A)T=A$.
\end{lemma}

\begin{lemma}
    \label{still invertible}
    Let $T$ be an integral domain with quotient field $K$, $R$ a subring with identity, and $I=(R:T)$. If $J$ is an invertible integral ideal of $R$, then $JT$ is an invertible integral ideal of $T$. If $A$ is an invertible integral ideal of $T$ which is comaximal to $I$, then $R\cap A$ is an invertible integral ideal of $R$ which is comaximal to $I$.
\end{lemma}

\begin{theorem}
    \label{exact sequence}
    Let $T$ be an integral domain, $R$ a subring with identity, and $I=(R:T)$. Consider the following sequence:
    $$1\xrightarrow{\iota} U(R)\xrightarrow{\phi} U(T)\oplus U(\sfrac{R}{I})\xrightarrow{\psi} U(\sfrac{T}{I})\xrightarrow{\sigma} \Cl(R)\xrightarrow{\tau}\Cl(T)\xrightarrow{\pi}1$$
    with maps $\phi$, $\psi$, and $\sigma$ defined as follows:
    \begin{alignat*}{2}
        &\phi:U(R)\to U(T)\oplus U(\sfrac{R}{I}),&&\quad\quad\phi(v)=(v,v^{-1}+I);\\
        &\psi:U(T)\oplus U(\sfrac{R}{I})\to U(\sfrac{T}{I}),&&\quad\quad\psi(u,r+I)=ur+I;\\
        &\sigma:U(\sfrac{T}{I})\to \Cl(R),&&\quad\quad\sigma(t+I)=[tT\cap R];\\
        &\tau:\Cl(R)\to\Cl(T),&&\quad\quad\tau([J])=[JT].
    \end{alignat*}
    Then $\phi$, $\psi$, $\sigma$, and $\tau$ are homomorphisms. Moreover, the sequence $(\iota,\phi,\psi,\sigma)$ is exact. The sequence $(\iota,\phi,\psi,\sigma,\tau)$ is exact if and only if every ideal class in $\ker(\tau)$ (i.e., any class $[J]\in\Cl(R)$ such that $JT$ is principal) contains an ideal class representative $J$ which is an integral ideal of $R$ comaximal to $I$. If, in addition to the previous condition, every ideal class in $\Cl(T)$ contains an ideal class representative $J$ which is an integral ideal of $T$ comaximal to $I$, then the sequence $(\iota,\phi,\psi,\sigma,\tau,\pi)$ is exact.
\end{theorem}

\begin{proof}
    To start, we will show that the maps $\phi$, $\psi$, $\sigma$, and $\tau$ are (well-defined) homomorphisms. Then, we will show exactness.

    First, note that for any $v\in U(R)$, $v^{-1}\in R\subseteq T$. Then $v\in U(T)$ as well. Furthermore, $(v+I)(v^{-1}+I)=1+I$, so $v^{-1}+I\in U(\sfrac{R}{I})$. Then $\phi$ is a well-defined function. Moreover, note that for any $v_1,v_2\in U(R)$,
    $$\phi(v_1v_2)=(v_1v_2,(v_1v_2)^{-1}+I)=(v_1,v_1^{-1}+I)(v_2,v_2^{-1}+I)=\phi(v_1)\cdot \phi(v_2).$$
    Thus, $\phi$ is a homomorphism.

    Now let $u\in U(T)$ and $r+I\in U(\sfrac{R}{I})$. Then for some $s\in R$, $(r+I)(s+I)=1+I$. For this choice of $s$, note that $(ur+I)(u^{-1}s+I)=1+I$ in $\sfrac{T}{I}$, so $ur+I\in U(\sfrac{T}{I})$. Moreover, if $u\in U(T)$ and $r_1+I=r_2+I\in U(\sfrac{R}{I})$, then $r_1-r_2\in I$. Therefore, $ur_1-ur_2=u(r_1-r_2)\in I$, so $ur_1+I=ur_2+I$. Thus, $\psi$ is a well-defined function. To see that $\psi$ is a homomorphism, note that for $u_1,u_2\in U(T)$ and $r_1+I,r_2+I\in U(\sfrac{R}{I})$,
    $$\psi(u_1u_2,r_1r_2+I)=u_1u_2r_1r_2+I=(u_1r_1+I)(u_2r_2+I)=\psi(u_1,r_1+I)\cdot\psi(u_2,r_2+I).$$

    Now let $t+I\in U(\sfrac{T}{I})$; then $J:=tT\cap R$ is an ideal of $R$. Since $tT$ is a principal (and thus invertible) ideal of $T$ comaximal to $I$, then Lemma \ref{still invertible} tells us that $J$ is an invertible ideal of $R$ comaximal to $I$. Then $[J]=[tT\cap R]\in \Cl(R)$. Letting $t_1+I=t_2+I\in U(\sfrac{T}{I})$, we note that $t_1-t_2\in I$. Then for any $s\in T$ such that $t_1s\in R$, $t_2s=t_1s+(t_2-t_1)s\in R$. By symmetry of $t_1$ and $t_2$, this tells us that for any $s\in T$, $t_1s\in R$ if and only if $t_2s\in R$. Then $t_1T\cap R=\frac{t_1}{t_2}(t_2T\cap R)$, so $[t_1T\cap R]=[t_2T\cap R]$ in $\Cl(R)$. Thus, $\sigma$ is a well-defined function.

    To show that $\sigma$ is a homomorphism, let $t_1+I,t_2+I\in U(\sfrac{T}{I})$. Trivially, $(t_1T\cap R)(t_2T\cap R)\subseteq t_1t_2T\cap R$. To show the reverse inclusion, let $\alpha\in T$ such that $t_1t_2\alpha\in R$. Since $t_1$ and $t_2$ are units modulo $I$, there exist $s_1,s_2\in T$ and $\beta_1,\beta_2\in I$ such that $t_1s_1+\beta_1=t_2s_2+\beta_2=1$ (in particular, $t_1s_1,t_2s_2\in R$). Then:
    \begin{align*}
        t_1t_2\alpha&=t_1t_2\alpha(t_1s_1+\beta_1)(t_2s_2+\beta_2)\\&=[t_1s_1][t_2s_2\cdot t_1t_2\alpha]+[t_1s_1][t_2(t_1\alpha\beta_2)]+[t_1(t_2\alpha\beta_1)][t_2s_2]+[t_1(\alpha\beta_1)][t_2\beta_2].
    \end{align*}
    Since each addend on the right-hand side of this identity is an element of $(t_1T\cap R)(t_2T\cap R)$, then $t_1t_2\alpha\in (t_1T\cap R)(t_2T\cap R)$. Thus, $$\sigma(t_1t_2+I)=[t_1t_2T\cap R]=[t_1T\cap R][t_2T\cap R]=\sigma(t_1+I)\cdot \sigma(t_2+I),$$
    so $\sigma$ is a homomorphism.

    Now let $[J]\in \Cl(R)$; without loss of generality, we may assume that $J$ is an integral ideal of $R$. By Lemma \ref{still invertible}, $JT$ is an invertible ideal of $T$, so $[JT]\in \Cl(T)$. Furthermore, if $[J_1]=[J_2]$, i.e., if there exist nonzero $r_1,r_2\in R$ such that $J_1=\frac{r_1}{r_2}J_2$, then $J_1T=\frac{r_1}{r_2}J_2T$. Thus, $\tau([J_1])=[J_1T]=[J_2T]=\tau([J_2])$, so $\tau$ is a well-defined function. To see that $\tau$ is in fact a homomorphism, note that for any $[J_1],[J_2]\in\Cl(R)$, $$\tau([J_1J_2])=[J_1J_2T]=[J_1T][J_2T]=\tau([J_1])\cdot\tau([J_2]).$$

    We now need to show exactness of the sequence $(\iota,\phi,\psi,\sigma)$. To do so, we need to show the following:
    \begin{align*}
        \{1\}&=\ker(\phi);\\
        \phi(U(R))&=\ker(\psi);\\
        \psi(U(T)\oplus U(\sfrac{R}{I}))&=\ker(\sigma).
    \end{align*}

    First, we will show that $\ker(\phi)$ is trivial; that is, that $\phi$ is injective. Let $u\in U(R)$ such that $\phi(u)=(u,u^{-1}+I)=(1,1+I)$. Then trivially $u=1$, so $\ker(\phi)=\{1\}$.

    We now want to show that $\phi(U(R))=\ker(\psi)$. First, let $v\in U(R)$ and note that $\psi(\phi(v))=\psi(v,v^{-1}+I)=vv^{-1}+I=1+I$. Then $\phi(v)\in \ker(\psi)$, so $\phi(U(R))\subseteq\ker(\psi)$. For the reverse inclusion, let $(t,r+I)\in \ker(\psi)$, i.e., $t\in U(T)$ and $r+I\in U(\sfrac{R}{I})$ such that $\psi(t,r+I)=tr+I=1+I$. Since $r+I\in U(\sfrac{R}{I})$, there exist $s\in R$ and $\beta_1\in I$ such that $sr+\beta_1=1$; since $tr+I=1+I$, there exists $\beta_2\in I$ such that $tr+\beta_2=1$. Then $t-s=(t-s)(sr+\beta_1)=(tr-sr)s+(t-s)\beta_1=(\beta_1-\beta_2)s+(t-s)\beta_1\in I$. Since $s\in R$ and $t-s\in I\subseteq R$, this means that $t\in R$. Moreover, $tr+\beta_2=1$ and $t^{-1}\in T$, so $t^{-1}=r+t^{-1}\beta_2\in R$; then $t\in U(R)$. Furthermore, $t^{-1}+I=(t+I)^{-1}=r+I$, so $t\in U(R)$ such that $(t,r+I)=(t,t^{-1}+I)=\phi(t)\in \phi(U(R))$. Thus, $\ker(\psi)\subseteq \phi(U(R))$.

    We will now show that $\psi(U(T)\oplus U(\sfrac{R}{I}))=\ker(\sigma)$. First, let $u\in U(T)$ and $r+I\in U(\sfrac{R}{I})$; then $\sigma(\psi(u,r+I))=\sigma(ur+I)=[urT\cap R]=[rT\cap R]$. Since $rT=(rR)T$, with $rR$ comaximal to $I$, then Lemma \ref{going up or going down} tells us that $(rR)T\cap R=rR$. Then $\sigma(\psi(u,r+I))$ is the principal (identity) class in $\Cl(R)$, so $\psi(U(T)\oplus U(\sfrac{R}{I}))\subseteq \ker(\sigma)$. To show the reverse inclusion, suppose that $t+I\in \ker(\sigma)$, i.e., $t+I\in U(\sfrac{T}{I})$ and $tT\cap R$ is a principal ideal in $R$. Then there exists $r\in R$ such that $tT\cap R=rR$. Since $t+I\in U(\sfrac{T}{I})$, there exists some $s\in T$ and $\beta\in I$ such that $st+\beta=1$. Since $st=1-\beta\in tT\cap R$, then $st=r\gamma$ for some $\gamma\in R$, so $r\gamma+\beta=1$. Then $r+I\in U(\sfrac{R}{I})$. Since $t$ is comaximal to $I$, Lemma \ref{going up or going down} tells us that $tT=(tT\cap R)T=(rR)T=rT$. Then $t$ and $r$ are associates in $T$, i.e., there exists $u\in U(T)$ such that $t=ur$, so $t+I=ur+I=\psi(u,r+I)\in \psi(U(T)\oplus U(\sfrac{R}{I}))$. Therefore, $\ker(\sigma)\subseteq \psi(U(T)\oplus U(\sfrac{R}{I}))$, so we conclude that $(\iota,\phi,\psi,\sigma)$ is an exact sequence.

    Now assume that every ideal class in $\ker(\tau)$ contains an ideal class representative $J$ which is an integral ideal of $R$ comaximal to $I$. Under this condition, we need to show that $\sigma(U(\sfrac{T}{I}))=\ker(\tau)$. Note that for any $t+I\in U(\sfrac{T}{I})$, the ideal $tT$ is comaximal to $I$. Then Lemma \ref{going up or going down} tells us that $\tau(\sigma(t+I))=\tau([tT\cap R])=[(tT\cap R)T]=[tT]$, the principal (identity) class in $\Cl(T)$. Then $\sigma(U(\sfrac{T}{I}))\subseteq\ker(\tau)$. To show the reverse inclusion, suppose that $[J]\in \ker(\tau)$, i.e., $JT=tT$ is principal. By our assumption, we can choose $J$ to be an integral ideal of $R$ comaximal to $I$. Then $R=J+I$, so $T=RT=(J+I)T=JT+I$. Then $JT=tT$ is comaximal to $I$, so $t+I\in U(\sfrac{T}{I})$. Then by Lemma \ref{going up or going down}, $tT\cap R=JT\cap R=J$, so $\sigma(t+I)=[tT\cap R]=[J]$. Thus, $[J]\in \sigma(U(\sfrac{T}{I}))$, so $\sigma(U(\sfrac{T}{I}))=\ker(\tau)$. Therefore, if every ideal class in $\ker(\tau)$ contains an ideal class representative $J$ which is an integral ideal of $R$ comaximal to $I$, then the sequence $(\iota,\phi,\psi,\sigma,\tau)$ is exact.

    On the other hand, assume that the sequence $(\iota,\phi,\psi,\sigma,\tau)$ is exact. In particular, this means that $\sigma(U(\sfrac{T}{I}))=\ker(\tau)$. Then every ideal class in $\ker(\tau)$ can be represented as $\sigma(t+I)=[tT\cap R]$ for some $t+I\in U(\sfrac{T}{I})$. Clearly $tT\cap R$ is an integral ideal of $R$; since $t+I\in U(\sfrac{T}{I})$, then $tT$ is comaximal to $I$ in $T$. Then by Lemma \ref{still invertible}, $tT\cap R$ is comaximal to $I$ in $R$. Thus, every ideal class in $\ker(\tau)$ contains an ideal class representative $J$ which is an integral ideal of $R$ comaximal to $I$.

    Finally, assume that, in addition to the previous condition, every ideal class in $\Cl(T)$ contains an ideal class representative $J$ which is an integral ideal of $T$ comaximal to $I$. In order to show that the sequence $(\iota,\phi,\psi,\sigma,\tau,\pi)$ is exact, we only need to show that $\tau$ is surjective. To do so, let $[A]\in\Cl(T)$. By the assumption, we can without loss of generality assume that $A$ is an integral ideal of $T$ comaximal to $I$. Then by Lemma \ref{still invertible}, $R\cap A$ is an invertible integral ideal of $R$ comaximal to $I$; that is, $[R\cap A]\in \Cl(R)$. Moreover, by Lemma \ref{going up or going down}, $\tau([R\cap A])=[(R\cap A)T]=[A]$. Then $\tau$ is surjective, so the sequence $(\iota,\phi,\psi,\sigma,\tau,\pi)$ is exact.
\end{proof}

\end{document}
